\section{Carré d'un nombre et racine carrée}
\subsection{Carré d'un nombre}

\begin{definitionbox}
Le \bb{carré} d'un nombre est le produit de ce nombre par lui même.
$$x^2 = x \times x$$
\end{definitionbox}

\begin{exemplebox}
Le carré de 7 est : $ 7^2 = \answer{7\times 7} = \answer{49}$\\
Le carré de 3.4 est : $ (3.4)^2 = \answer{3.4 \times 3.4} = \answer{11.56}$\\
    Le carré de -5 est : $ (-5)^2 = \answer{(-5) \times (-5)} = \answer{+25}$
\end{exemplebox}

\begin{propbox}
    Le carré d'un nombre est \bb{toujours positif}. 
\end{propbox}

\subsection{Racine carrée d'un nombre}


\begin{definitionbox}
    La racine carrée d'un nombre \bb{positif} est l'unique nombre positif qui \bb{mis au carré}
    est égale à ce nombre. 
    \newline
    \newline
    \bb{Notation} : Si $a$ est un nombre positif, la racine carrée
    de $a$ sera notée $\sqrt{a}$.
\end{definitionbox}

\begin{exemplebox}
    \begin{itemize}
        \item La racine carrée de 49 est : $ \sqrt{49} =\sqrt{7\times 7}=
            \answer{\sqrt{7^2} = 7}$
        \item La racine carrée de 144 est : $ \sqrt{144} = \answer{\sqrt{12
            \times 12} = \sqrt{12^2} = 12}$
        \item La racine carrée de 2 est : $ \sqrt{2} \approx 1,141213562$
    \end{itemize}
\end{exemplebox}

\begin{definitionbox}
    On dit qu'un nombre est un \bb{carré parfait} si \answer{sa racine carrée est
    un nombre entier.}
\end{definitionbox}

\begin{remarquebox}
    Lorsque la racine carrée d'un nombre est un nombre décimal avec de nombreux
    chiffres après la virgules, on utilise le symbole $\approx$ pour donner une
    \bb{valeur approchée} de la racine carrée.
    \newline
    \newline
    Par exemple on peut donner une approximation au centième de $\sqrt{11}$ :
    $$\sqrt{11} = 3.316625... \approx \answer{3.32} $$
    Un approximation au millième serait $\answer{3.317}$. \hl[RougeClair]{La
    valeur exacte de $\sqrt{11}$ est $\sqrt{11}$}.
\end{remarquebox}

\begin{aretenir}
    Il est important de retenir les \bb{carrés parfaits} de 1 à 15. Ces nombres
    nous permettent de vérifier rapidement la vraissemblance de nos calculs
    lorsqu'on manipule les carrés et les racines carrées.
    \begin{center}
        \begin{tabular}{|c||c|c|c|c|c|c|c|c|c|c|c|c|c|c|c|}
        \hline
            $n$ & 1 & 2 & 3 & 4  & 5  & 6  & 7  & 8  & 9  & 10  & 11  & 12  & 13  & 14  & 15  \\ \hline
            $n^2$ & 1 & 4 & 9 & 16 & 25 & 36 & 49 & 64 & 81 & 100 & 121 & 144 & 169 & 196 & 225 \\ \hline
        \end{tabular}
    \end{center}
\end{aretenir}

\begin{propbox}[\it{Carré de la racine et racine du carré}]
    Pour tout nombre $a$ \bb{positif} on a : 
    $$(\sqrt{a})^2=\sqrt{a^2}=a$$
\end{propbox}

\begin{methodebox}
    Soit un $a$ un nombre. Encadrer $\sqrt{a}$ par deux entiers consécutifs
    c'est trouver un entier $n$ tel que : 
    $$n\le \sqrt{a}\le n+1 $$
\end{methodebox}

\begin{exemplebox}
    Pour encadrer $\sqrt{20}$ on cherche deux entiers qui mis au
    carrés encadrent $\sqrt{20}$. Comme $4^2 = 16$ et $5^2=25$, on a bien
    $$4^2 \le 20 \le 5^2$$
    Ce qui nous permet d'affirmer que :
    $$4 \le \sqrt{20} \le 5$$
\end{exemplebox}


